\chapter{Correspondencia con la rúbrica de evaluación}\label{anx:rubrica}

\begin{table}[H]
\centering\small
\begin{tabularx}{\textwidth}{>{\raggedright\arraybackslash}p{4.4cm}r>{\raggedright\arraybackslash}X}
\toprule
Criterio de la rúbrica & Puntos & Dónde se desarrolla \\
\midrule
Justificación científica & 2 & Sección~\ref{sec:justif} \\
Hipótesis, objetivos & 2 & Secciones~\ref{sec:pregunta}--\ref{sec:hipotesis} y tabla~\ref{tab:hipotesis} \\
Revisión de antecedentes teóricos y experimentales & 2 & Capítulo~\ref{cap:estado} (secciones~\ref{sec:ant_teo} y \ref{sec:ant_exp}) y anexo~\ref{anx:busqueda} \\
Presentación de resultados teóricos y/o experimentales & 4 & Teóricos: capítulos~\ref{cap:marco}--\ref{cap:diseno}. Experimentales: capítulo~\ref{cap:resultados} (tabla~\ref{tab:resumen_res}) \\
Discusión de resultados & 2 & Capítulo~\ref{cap:discusion} \\
Conclusiones & 2 & Capítulo~\ref{cap:conclusiones} (tabla~\ref{tab:correspondencia}) \\
Bibliografía & 1 & 44 referencias en APA 7, con DOI verificados en Crossref y todas citadas en el texto \\
\bottomrule
\end{tabularx}
\end{table}

\chapter{Registro de la búsqueda bibliográfica}\label{anx:busqueda}

\textbf{Fuentes:} Google Scholar, RePEc, arXiv, NBER, Crossref; repositorios peruanos (ALICIA-CONCYTEC, UNI, UPeU, ESAN, Universidad de Lima, Universidad del Pacífico, IEP), BCRP y normativa municipal. Los metadatos se verificaron el 21 y el 25 de octubre de 2026.

\textbf{Cadenas principales:}
\begin{itemize}[itemsep=0pt]
\item \path{("double machine learning" OR "causal forest") AND (housing OR rental OR hedonic)}
\item \path{"precio alquiler departamentos Lima" AND (hedónico OR "machine learning" OR estacionamiento OR cochera)}
\item \path{("inferencia causal" OR "diferencias en diferencias") AND vivienda AND (Lima OR Perú)}
\item \path{("bundled parking" OR "parking premium") AND rent}
\item \path{("large language model" OR GPT) AND ("real estate" OR housing) AND listing}
\end{itemize}

\textbf{Criterios:} se incluyeron estudios que estiman precios implícitos, predicen precios o identifican efectos de atributos de la vivienda, y los trabajos metodológicos usados directamente. Se excluyeron los estudios sin metadatos verificables y las versiones duplicadas, en cuyo caso se cita la versión publicada.

\chapter{Reproducibilidad y trazabilidad}\label{anx:inventario}

\begin{table}[H]
\centering\footnotesize
\begin{tabularx}{\textwidth}{>{\raggedright\arraybackslash}p{5cm}X}
\toprule
Artefacto (\path{notebooks/}) & Contenido \\
\midrule
\path{notebook2_auditoria} & R1: limpieza, duplicados, ROOFTOP; genera \path{csv/data_causal_base.csv} \\
\path{notebook2b_extraccion_ia} + \path{extraccion_ia.py} & R8: extracción de la cochera (regex + OpenAI, con caché) \\
\path{notebook3b_efecto_cochera} & R2--R5, R7: resultado principal \\
\path{notebook3_estimacion_causal} & Análisis secundario (baños) \\
\path{notebook4_semisintetico} + \path{semisintetico.py} & R6: validación semisintética \\
\path{notebook_resumen} & EDA resumido y recorrido de toda la tesis \\
\path{causal_utils.py} & Estimadores DML (PLR, IRM), GATE, sensibilidad, Moran \\
\path{informe/tablas_manuscrito.py} & Convierte los CSV de resultados en las tablas del informe \\
\bottomrule
\end{tabularx}
\end{table}

\textbf{Orden de ejecución:} \path{notebook2} $\rightarrow$ \path{2b} $\rightarrow$ \path{3b} $\rightarrow$ \path{3} $\rightarrow$ \path{4} $\rightarrow$ \path{notebook_resumen} $\rightarrow$ \path{tablas_manuscrito.py} $\rightarrow$ compilación. La semilla es \path{SEED = 42}. Una segunda ejecución completa reprodujo exactamente las 29 tablas y los resúmenes.

\textbf{Software:} Python 3.13, pandas 3.0, scikit-learn 1.9 \citep{pedregosa2011}, statsmodels 0.15 y \texttt{gpt-4.1-mini-2025-04-14}, con el prompt v2.

\textbf{Privacidad:} los datos (\path{data/}, \path{csv/}) y la clave (\path{.env}) no se versionan, y la base analítica no contiene descripciones, direcciones ni teléfonos.
